\section{Experimental Setup}
\label{sec:setup}

\subsection{Dataset: Defects4J v2.0}

Experiments use \textit{Defects4J}~v2.0, which provides over 830 active bugs
across 17 open-source Java projects \citep{just2014defects4j}.
Table~\ref{tab:projects} lists the projects with approximate per-project active
counts. From this benchmark we pre-collected a shared evidence corpus of 854
buggy and 854 fixed \texttt{context.json} payloads spanning all 17 projects; every
study reuses this corpus, so no manifest triggers fresh \textit{Defects4J}
checkout or test execution.

\begin{table}[H]
\centering
\caption{\textit{Defects4J}~v2.0 projects and approximate active-bug counts.
         Per-project counts are approximate and version-dependent; the shared
         evidence corpus covers 854 bugs across these projects.}
\label{tab:projects}
\setlength{\tabcolsep}{6pt}
\begin{tabular}{@{}llr@{}}
\toprule
\thead{Project} & \thead{Description} & \thead{Bugs} \\
\midrule
Chart           & JFreeChart plotting library     & 26  \\
Cli             & Apache Commons CLI              & 40  \\
Closure         & Google Closure Compiler         & 174 \\
Codec           & Apache Commons Codec            & 18  \\
Collections     & Apache Commons Collections      & 4   \\
Compress        & Apache Commons Compress         & 47  \\
Csv             & Apache Commons CSV              & 16  \\
Gson            & Google Gson                     & 18  \\
JacksonCore     & Jackson JSON core               & 26  \\
JacksonDatabind & Jackson data binding            & 112 \\
JacksonXml      & Jackson XML                     & 6   \\
Jsoup           & HTML parser                     & 93  \\
JxPath          & Apache JXPath                   & 22  \\
Lang            & Apache Commons Lang             & 64  \\
Math            & Apache Commons Math             & 106 \\
Mockito         & Mocking framework               & 38  \\
Time            & Joda-Time                       & 26  \\
\midrule
\textbf{Total}  &                                 & \textbf{836} \\
\bottomrule
\end{tabular}
\end{table}

\subsection{Context File Structure}

Section~\ref{sec:collection} describes how each bug's case file is built; here
we list exactly what it contains. Every \texttt{context.json} holds: the bug ID,
project name, and report URL; failing-test summaries, exception messages, and
stack-trace frames; production-code and failing-test snippets; optional
Cobertura coverage summaries; hidden-oracle fields (\texttt{classes.modified} and
the optional fix diff); and the buggy/fixed revision-pair metadata.

\subsection{Execution Environment}

Classifications were executed against the Google Gemini API (model
\texttt{gemini-3.1-flash-lite-preview}) with schema-enforced structured output.
\textit{Defects4J} tooling is invoked through WSL for the collection stage, but
because evidence is reused from the pre-collected corpus, the reported runs
involve only LLM calls. Commands are scripted and dependency versions pinned for
reproducibility, and each run records context-reuse counts and per-classification
turn/call counts.

\subsection{Evidence Collection Configuration}

Stack-frame filtering retains only project source frames, discarding JUnit,
TestNG, Hamcrest, JDK-internal, and build-tool frames. Production snippets capture
$\pm$12 lines around each suspicious frame; test snippets capture $\pm$18 lines
and are limited to the first three failing tests per bug. Coverage analysis is
optional and focuses on suspicious classes parsed from Cobertura XML.

\subsection{Condition Configuration}

Each research question fixes the taxonomy and strategy it needs
(Table~\ref{tab:rqmap}). The default \textit{scientific-open} configuration
supplies RQ1, RQ3, and RQ5; RQ2 additionally runs the \textit{scientific-closed}
pass; and RQ4 runs the \textit{zero-free} and \textit{few-open} rungs. The
few-closed condition is valid but is not required by any current question and is
omitted. All conditions draw on the same evidence corpus and the same prompt
components described in Section~\ref{sec:strategies}. Table~\ref{tab:rqmap}
lays out, for every research question, which metrics answer it and which
analysis command computes them. 

\begin{table}[H]
\centering
\caption{Research questions, the metrics they use, and the analysis command that
         computes each.}
\label{tab:rqmap}
\setlength{\tabcolsep}{5pt}
\begin{tabularx}{\columnwidth}{@{}lXl@{}}
\toprule
\thead{RQ} & \thead{Metric(s)} & \thead{Command} \\
\midrule
RQ1 & ODC type distribution; per-project variation     & \texttt{study-drift}  \\
RQ2 & Escape rate; coverage rate; closed$\leftrightarrow$open shift-$\kappa$
                                                        & \texttt{study-escape} \\
RQ3 & Four-tier: strict / top-2 / family / $\kappa$     & \texttt{study-drift}  \\
RQ4 & Unique labels; entropy; ODC coverage; vocab.\ reduction
                                                        & \texttt{study-ladder} \\
RQ5 & Pre/post drift transitions; per-project $\kappa$  & \texttt{study-drift}  \\
\bottomrule
\end{tabularx}
\end{table}

\subsection{Study Plan}

Two development-validation runs (a six-bug pilot and a larger balanced sample)
confirmed that every study command produces correctly tagged, well-formed metrics
end to end. These runs are not intended to yield the paper's quantitative
findings. The reported results (Section~\ref{sec:results}) are populated by a
\emph{confirmatory} run: a project-balanced manifest of 68--100 bugs
(\texttt{study-plan} with a per-project minimum of four to five, sufficient to
make per-project $\kappa$ non-degenerate everywhere), classified under the four
conditions the five questions require.

A multi-fault consideration bounds interpretation. Callaghan
\citep{callaghan2024mf} reports that real program versions contain an average of
about 9.2 co-existing faults, so symptom evidence from one fault can contaminate
the classification of another. The pipeline can query the \texttt{defects4j-mf}
dataset to flag such cases during analysis.

%% =============================================================================
